\pdfminorversion=7
\documentclass[11pt,a4paper]{article}
\usepackage[T1]{fontenc}
\usepackage[utf8]{inputenc}
\usepackage{lmodern,amsmath,amssymb,amsthm,mathtools}
\usepackage[margin=25mm,headheight=14pt]{geometry}
\usepackage{microtype,booktabs,tabularx,array,longtable,enumitem,fancyhdr,xurl}
\usepackage[hidelinks]{hyperref}
\newtheorem{theorem}{Theorem}
\newtheorem{lemma}[theorem]{Lemma}
\newtheorem{proposition}[theorem]{Proposition}
\newtheorem{corollary}[theorem]{Corollary}
\theoremstyle{definition}\newtheorem{definition}[theorem]{Definition}
\theoremstyle{remark}\newtheorem{remark}[theorem]{Remark}
\newcommand{\HH}{\mathsf H}
\newcommand{\Rea}{\operatorname{Re}}
\newcommand{\Ima}{\operatorname{Im}}
\newcommand{\Log}{\operatorname{Log}}
\newcommand{\dd}{\mathrm d}
\newcommand{\F}{\mathbb F}
\newcommand{\floor}[1]{\left\lfloor #1\right\rfloor}
\newcommand{\ceil}[1]{\left\lceil #1\right\rceil}
\setcounter{secnumdepth}{2}
\setlength{\parindent}{1.2em}
\setlength{\parskip}{3pt plus 1pt}
\setlength{\emergencystretch}{3em}
\setlist{nosep,leftmargin=1.5em}
\pagestyle{fancy}\fancyhf{}
\fancyhead[L]{\small Conditional bound by translated probes}
\fancyhead[R]{\small D. J. Raymond}
\fancyfoot[C]{\small\thepage}
\renewcommand{\headrulewidth}{0.3pt}
\hypersetup{pdftitle={Retuned translated probes and the conditional bound Lambda < 0.116398 for the de Bruijn-Newman constant},pdfauthor={Dale James Raymond},pdfsubject={First-deposit version 1.0; conditional computer-assisted theorem},pdfkeywords={de Bruijn-Newman constant, translated probes, matched logarithmic coefficients, interval arithmetic}}
\title{Retuned translated probes and the conditional bound\\
$\Lambda<0.116398$ for the de Bruijn--Newman constant}
\author{Dale James Raymond\\[3pt]
\small ORCID: \href{https://orcid.org/0009-0001-1682-8529}{0009-0001-1682-8529}\\
\small \href{mailto:dale_raymond@icloud.com}{\nolinkurl{dale_raymond@icloud.com}}}
\date{10 October 2026 --- Version 1.0}
\begin{document}\maketitle
\begin{abstract}
Assume that the Riemann zeta function has no nontrivial zeros in the open
half-plane $\Rea s>7/8$. A computer-assisted argument gives
\[
\Lambda\le\frac{116397726408291}{10^{15}}
=0.116397726408291<0.116398
\]
for the de Bruijn--Newman constant in the Polymath normalisation.
The proof combines a retuned probe-height schedule, matched finite
logarithmic coefficients and exact translated-kernel inequalities.
The source calculation preserves common multiplicative phases and cancels
shared convolution terms before bounding the remaining terms. A family of
44 rational kernel certificates then adapts the comparison to the smaller
height-ratio domains along the descent. Together these estimates improve
the complete contraction rate, including the entire distinguished-pair
subtraction. The physical evidence covers 608 closed time cells, six
frequency classes including an unbounded final class, and all eleven
translated class pairs. The retuned physical enclosures were produced in
two interval-arithmetic constructions. Two 14,899-step rational histories
reproduce byte-for-byte and pass a separately implemented auditor.
The paper gives the global first-contact argument, the corrected Gaussian
remainder and the continuous-domain kernel proof; the analytic supplements,
editable sources and executable evidence accompany it. The $7/8$ assertion
is an explicit hypothesis, not a conclusion of the numerical checks.
\end{abstract}
\noindent\textbf{2020 Mathematics Subject Classification.} Primary 11M26; secondary 65G30.\\
\textbf{Keywords.} de Bruijn--Newman constant; translated probes; matched logarithmic
coefficients; exact polynomial certificates; interval arithmetic.

\section{Introduction}\label{sec:introduction}
The de Bruijn--Newman constant $\Lambda$ is the threshold at which all zeros
of the heat-deformed Riemann xi function are real. Heat-flow estimates can
combine a time-zero strip, finite localisation walls and
logarithmic-derivative bounds away from the zeros to give a quantitative
upper bound \cite{polymath,pt}. This paper strengthens the translated-probe
construction of \cite{raymond149} under an explicit narrower initial strip.

For $1/2<\alpha\le1$, let $\HH_\alpha$ mean that every nontrivial zeta zero
$\rho$ satisfies $\Rea\rho\le\alpha$. Zeros on the boundary are allowed.
\begin{theorem}[Conditional upper bound]\label{thm:main}
Assume $\HH_{7/8}$. In the normalisation of Section~\ref{sec:norm},
\begin{equation}\label{eq:bound}
\boxed{\Lambda\le\frac{116397726408291}{10^{15}}
=0.116397726408291<0.116398.}
\end{equation}
\end{theorem}

\subsection{Optimising the complete contraction rate}
A useful probe must do more than give a large logarithmic-derivative lower
bound. Its height, its translation and its kernel constant also determine
how much zero contraction that bound supplies. The quantity to improve is
\[
 c=\frac{2\lambda}{h}\left(\ell-\frac Mh\right),
\]
where $\ell$ is the complete physical source floor and $M/h$ removes the
whole contribution of the distinguished conjugate pair. Optimising only
$\ell$ or only $\lambda$ can choose the wrong probe.

The first step makes the trade-off concrete. The earlier prepared
comparison used height $2.18$, a source floor about $6.3340$ and rate about
$10.3210$. The present comparison uses height $1.88$, a smaller source
floor about $6.0650$, and a larger rate about $11.1631$. The smaller source
floor is more useful after the height and own-pair terms are included.
These rounded values illustrate exact rational data, not sampled physical
measurements.

Three changes supply the improvement. First, the kernel family certifies
larger constants on smaller height-ratio domains. Second, the physical
estimates are regenerated at 274 changed probe heights. Third, the
matched-logarithmic calculation retains the shared finite convolution
before estimating the remainder. The final schedule selects one complete
certified source row in each time cell and is frozen before the exact
histories are generated. No source bound is moved to an uncertified height.

\begin{center}\small
\begin{tabularx}{\textwidth}{@{}Xr@{}}\toprule
Construction in the retained development history & Exact primary ceiling\\\midrule
Earlier undeposited comparison & 0.118068908899880\\
New kernels; unchanged physical source & 0.117794794642445\\
Retuned heights; basic joint source & 0.116787767525899\\
Retuned heights; matched logarithmic source & \textbf{0.116397726408291}\\\bottomrule
\end{tabularx}
\end{center}
The total reduction is $0.001671182491589$, approximately $1.41543\%$ of
the earlier draft ceiling. All four comparisons use the same strict seed
$9/16+10^{-6}$ and 32 substeps per source cell. The gain is therefore not
from tightening the seed or increasing the time subdivision. This is a
comparison of completed constructions, not a global-optimality claim.

\subsection{Context, assumptions and scope}
The $7/8$ assertion appears in the supplied OpenAI manuscript of
30 September 2026 \cite{openai78}. Its 5 October companion states the
different $11/12$ half-plane \cite{openai1112}. The pinned formal-source
declaration is identified in \cite{openaiLean}; no execution of that
formalisation is claimed here. The mathematical assumption used in this
paper is exactly $\HH_{7/8}$, regardless of the external assertion's
publication or review status.

Polymath's heat-flow framework and the finite-RH result of Platt and
Trudgian are analytic inputs \cite{polymath,pt}. The author's earlier
public paper reports $\Lambda<0.149396$ from the classical initial strip
\cite{raymond149}. That result uses a different initial hypothesis from the one here.
The subsequent preprints discussed there, including Gomila and Gordon
\cite{gomila,gordon}, and Planat's collective-contraction framework
\cite{planat}, provide context. No numerical certificate from those three
works is imported into Theorem~\ref{thm:main}. 

The companion \emph{Further Consequences} records a separate implication
from $\HH_{34/43}$ to $\Lambda<0.09$, its still-unproved covariance supplier,
and a restricted limitation of the older source-table method. None is a
premise of Theorem~\ref{thm:main}.

\subsection{Reading the proof}
Sections~\ref{sec:norm}--\ref{sec:matched} give the initial strip,
localisation and finite logarithmic source argument.
Section~\ref{sec:retuned} specifies the retuned source and its exact
selection. Sections~\ref{sec:kernel}--\ref{sec:comparison} prove kernel
domination and global comparison. Section~\ref{sec:replay} identifies the
exact endpoint and verification boundary. The two appendices give the
complete kernel-constant table and the Gaussian correction.
The supplied \emph{Matched Source Supplement} expands the finite-source
algebra; \emph{Prior Analytic Foundations} gives the physical approximation
and localisation framework, with its Gaussian bound replaced explicitly
by Appendix~\ref{app:gaussian} below.

\section{Normalisation, initial separation and localisation}\label{sec:norm}
Use the normalisation of \cite{polymath}:
\begin{align}
H_t(z)&=\int_0^\infty e^{tu^2}\Phi(u)\cos(zu)\,\dd u,\label{eq:H}\\
\Phi(u)&=\sum_{n\ge1}(2\pi^2n^4e^{9u}-3\pi n^2e^{5u})
 e^{-\pi n^2e^{4u}},\nonumber\\
H_0(z)&=\tfrac18\xi((1+iz)/2),\qquad
\partial_t H_t=-\partial_z^2H_t.\nonumber
\end{align}
Here $\xi(s)=\tfrac12s(s-1)\pi^{-s/2}\Gamma(s/2)\zeta(s)$.
The functions $H_t$ are real and even. The de Bruijn--Newman constant is the
threshold such that $H_t$ has only real zeros exactly when $t\ge\Lambda$.
De Bruijn's strip theorem, in \cite[Theorem 3.2]{polymath}, contracts a
zero strip of half-height $Y$ at time $t$ to half-height at most
$\sqrt{\max(Y^2-2\tau,0)}$ at time $t+\tau$.

\begin{lemma}[Time-zero strip]\label{lem:initial}
Under $\HH_{7/8}$ every zero of $H_0$ has $|\Ima z|\le3/4$.
The initial squared comparison value used here is
\begin{equation}\label{eq:q0}
q_0=\frac{562501}{1000000}>\frac9{16},\qquad
q_0-\frac9{16}=\frac1{1000000}.
\end{equation}
\end{lemma}
\begin{proof}
The functional equation reflects a nontrivial zero $\rho$ to $1-\rho$,
so $1/8\le\Rea\rho\le7/8$. Its corresponding $H_0$ height in
\eqref{eq:H} is $1-2\Rea\rho$. The strict seed inequality is exact.
\end{proof}
Ordinary contraction from this strip alone gives $\Lambda\le9/32$.
The physical estimates below supply the additional contraction needed for
the stated bound.

Fix
\[
X_0=6000000185827,\qquad N_0=690988.
\]
The function $B_t$ is the nonvanishing analytic normaliser defined in the
included Foundations, Part A, Section 4, with the same logarithm branches
and normalisation throughout. The finite-head wall is valid for $x=X_0$, $0\le t\le41/250$ and
$0\le y\le1$. With the corrected physical error allowance in
Appendix~\ref{app:gaussian}, the unchanged ideal-prefix certificates imply
\begin{equation}\label{eq:wall}
\left|\frac{H_t(X_0+iy)}{B_t(X_0+iy)}\right|
>\frac{79}{100}-\frac1{100}-\frac1{180}
=\frac{697}{900}>\frac34.
\end{equation}
Both complete wall constructions are included. The finite-RH theorem of
Platt--Trudgian verifies the exact ordinate $3000175332800$, which exceeds
$(X_0+1)/2$ by $175239886$ \cite[Theorem 1]{pt}. Its original zero
computation is cited, not repeated.

Initially every zero inside the finite head is real. The wall, positivity
of $H_t(iy)$, the enclosing classical strip, and local forward Hermite
splitting at a real multiple zero prevent entry or birth of a nonreal zero
in that head during the declared time range. Thus every prospective
nonreal contact is at $|\Rea z|>X_0$. These are the localisation arguments
of the included Foundations, Part A, Sections 2--3, with the safe wall
\eqref{eq:wall}. No later narrow-strip or finite-canopy result is imported.

\section{Uniform physical source bounds at all frequencies}\label{sec:physical}
Define, only after proving probe nonvanishing,
\[
S(x,h,t)=-\Ima\frac{H_t'(x+ih)}{H_t(x+ih)}.
\]
The selected source in \nolinkurl{evidence/SOURCE_SELECTED.json} prescribes rational
heights $h_i$ on the 608 closed cells
\[
I_i=[i/4000,(i+1)/4000],\qquad 0\le i<608.
\]
Its heights lie in $[69/100,47/25]$. The schedule is part of the theorem's finite data. Each comparison step
uses the lower bound certified at its prescribed height, without
interpolation between heights.
For $s=\tfrac12\log(x/(4\pi))-\log N_0$, the six frequency classes are
\[
[0,1/20],\ [1/20,3/20],\ [3/20,2/5],\ [2/5,1],\ [1,2],\ [2,\infty).
\]
They cover every actual $x\ge X_0-2$. The source activation bounds include
all native cutoffs, including future cutoffs in the unbounded final class.
The translated points $x\pm a$, with $x\ge X_0$ and $|a|<2$, belong to
one of the eleven class pairs $(j,j)$ or $(j,j+1)$. The original source proof establishes this same-or-adjacent-class
restriction and handles the joins between classes.

\subsection{Physical approximation and the labelled model}
Put $L=\log(x/(4\pi))$, $b_t(n)=\exp(t\log^2n/4)$ and
\[
s_U^0=(1+h)/2+tL/4-i(x/2+\pi t/8),\qquad
E_t(s)=\prod_{p\in\{2,3,5\}}(1-b_t(p)p^{-s}).
\]
Multiplying the \emph{finite} first and reflected branches by $E_t$ gives
the complete labelled model
\[
Q=1+\sum_{j>0}a_jz_j,\qquad
V=\sum_{j>0}\lambda_ja_jz_j,\qquad Q_h=-V/2,
\]
with $a_j,\lambda_j\ge0$, $|z_j|=1$ and all reflected and product-tail
labels retained. Certified masses satisfy
\[
\bar A\ge\sum a_j,\quad \bar B\ge\sum\lambda_ja_j,
\quad\bar A<4/5.
\]
The source first proves nonvanishing, then divides to obtain
\begin{equation}\label{eq:physical}
S>\ell+\frac{C-\Rea(V/Q)}2+
\left(E_\ell-\Rea(E_h/E)\right).
\end{equation}
The floor $\ell$ already includes subtraction of the physical errors.
The term $E_\ell$ is the source's phase-free upper cap for
$\Rea(E_h/E)$; their difference is a nonnegative phase deficit.

The all-real-abscissa Riemann--Siegel/Gaussian derivation and its frozen-native-
cutoff Cauchy disk are in the Foundations, Part A, Sections 4--5. The
actual disks obey $x\ge X_0-4$, $0.54\le y\le2.41$, and
$0\le t\le19/125$. The calculation includes the value error, the complete edge pair and the
derivative error, and proves denominator exclusion. Appendix~\ref{app:gaussian}
corrects the coarse Gaussian comparison on this domain and checks the complete
physical-error rows of the original and retuned source constructions. The argument at $h>1$ uses the enlarged-height
derivation rather than the restricted-height basic approximation theorem.
Cauchy's estimate is applied to the holomorphic error on the frozen-cutoff
disk, not to a nonholomorphic leading-model error.

\subsection{Retaining the common phase}
For $c_\pm\ge c_0\ge0$, $|z|=|u|=1$ and $\varepsilon\in\{1,-1\}$,
\begin{equation}\label{eq:phase}
\frac{c_+(1-\varepsilon\Rea(zu))+c_-(1-\varepsilon\Rea(z\bar u))}{2}
\ge c_0(1-|\Rea u|).
\end{equation}
Each factor $1-\varepsilon\Rea(zu^{\pm1})$ is nonnegative, so replacing
$c_\pm$ by $c_0$ gives a lower bound. The identity
$\Rea(zu)+\Rea(z\bar u)=2\Rea z\Rea u$ completes the estimate.
The original translated-probe refinement applies this only to identified
nonnegative cone-surplus terms and retains separate denominators at the
two points. It handles the constant/head pairs through index 1024 and the
first eight harmonics for the three finite Euler factors. Only the selected positive deficits contribute to the estimate; the
remaining deficits are discarded throughout the calculation.

\section{Matched finite logarithmic overlap}\label{sec:matched}
Matching the common finite-head terms gives an additional lower bound
for the paired source. The matched-source supplement, Sections 3--8,
contains the complete coefficient, support and numerical constructions. All logarithmic
identities below are formal coefficient identities; no convergent
positive-time Euler product is assumed.

\subsection{A formal positivity statement}
All algebra here is formal Dirichlet convolution:
$(f*g)(n)=\sum_{d\mid n}f(d)g(n/d)$. The function $\delta$ is its unit,
$\delta_p$ is one at $p$ and zero elsewhere, and $\log_*f$ is the formal
convolution logarithm of a series with $f(1)=1$. Write
\[
\mathcal Df(n)=(\log n)f(n),\qquad
c_t=b_t*\prod_{p\in\{2,3,5\}}(\delta-b_t(p)\delta_p),
\qquad \widetilde d_t=\mathcal D\log_*c_t.
\]
\begin{lemma}\label{lem:positive}
For every integer $n\ge1$ and $t\ge0$, the coefficients $c_t(n)$ and
$\widetilde d_t(n)$ and all their time derivatives are nonnegative.
Moreover $c_t(1)=1$ and $c_t(p)=\widetilde d_t(p)=0$ for $p=2,3,5$.
\end{lemma}
\begin{proof}
For $d_t=\mathcal D\log_*b_t$, the derivation rule and
$\partial_t b_t=\mathcal D^2b_t/4$ give
\[
\partial_t d_t(n)=\frac{\log^2n}{4}d_t(n)
+\frac{\log n}{4}\sum_{ab=n,\ a,b\ge2}d_t(a)d_t(b).
\]
The initial coefficients are the nonnegative von Mangoldt coefficients.
Induction on $n$ and the integrating-factor formula prove positivity of
all derivatives. At a selected prime power, put $\kappa=\log^2p/4$ and
$e_k=d_t(p^k)-(\log p)e^{k\kappa t}$. Then
\[
e_k'=k^2\kappa e_k+(k^2-k)\kappa(\log p)e^{k\kappa t}
+\frac{k\log p}{4}\sum_{j=1}^{k-1}d_t(p^j)d_t(p^{k-j}),\quad e_k(0)=0.
\]
Its forcing has all derivatives nonnegative. Elsewhere $\widetilde d=d$.
Finally the exact recurrence
\begin{equation}\label{eq:conv}
(\log n)c_t(n)=\sum_{d\mid n,\ d\ge2}\widetilde d_t(d)c_t(n/d)
\end{equation}
proves the assertion for $c$ by induction. The selected primes vanish
identically.
\end{proof}
For positive $t$, the infinite numerical heat-weighted Dirichlet series
converges at no fixed abscissa. The argument uses the preceding identities coefficientwise, and all
products evaluated below are finite.

\subsection{Cancelling shared terms before enclosure}
Let $\sigma=\Rea s_U^0$ and $K=8192$. In the complete first head through
$N_0$, the physical phases obey $z_dz_m=z_{dm}$. Define
\begin{align*}
T&=\sum_{2\le d\le K}\widetilde d_t(d)d^{-\sigma}z_d,
& D_K&=\sum_{2\le d\le K}\widetilde d_t(d)d^{-\sigma},\\
Z&=\sum_{2\le d\le K}\widetilde d_t(d)d^{-\sigma}
       \sum_{2\le m\le N_0/d}c_t(m)m^{-\sigma}.&&
\end{align*}
The matching uses only the first-head labels, where the multiplicative
relations hold. The reflected and product-tail labels remain separate.
By \eqref{eq:conv}, the shared mass between $QT$ and $V$ through the common
head is $D_K+Z$. Remove it from both sums before using the triangle inequality:
\begin{equation}\label{eq:resid}
|QT-V|\le\bar B+(1+\bar A)D_K-2(D_K+Z)
=\bar B-(1-\bar A)D_K-2Z.
\end{equation}
The right side bounds the mass of the unmatched labelled terms and is
therefore nonnegative. Using $|Q|\ge1-\bar A$ when dividing gives
\begin{equation}\label{eq:quotient}
\Rea(V/Q)\le
\frac{\bar B-2Z}{1-\bar A}-(D_K-\Rea T).
\end{equation}
The terms containing $D_K$ are combined algebraically before enclosure.
They do not receive separate upper and lower bounds that would lose the
cancellation.
Substitution in \eqref{eq:physical} gives
\begin{equation}\label{eq:newS}
S>F+\frac{D_K-\Rea T}{2}+
\left(E_\ell-\Rea(E_h/E)\right),
\qquad
F=\ell+\frac C2-\frac{\bar B-2Z_{\rm low}}{2(1-\bar A)}.
\end{equation}
The extra overlap is a positive finite sum, so it can be replaced by the
certified lower enclosure $Z_{\rm low}$. This estimate and the cone-surplus
estimate bound the same physical quantity. Either may be selected, but their
gains cannot be added.

\subsection{Paired logarithmic and grouped-head alternatives}
At $x\pm a$, the first-head phases satisfy
$z_d^\pm=z_d^c\exp(\pm ia\log d/2)$. From \eqref{eq:phase} and
Lemma~\ref{lem:positive}, a complete source-class pair $j\le k$ has the
following guaranteed logarithmic deficit:
\[
G_{k,a}=\sum_{2\le d\le K}\widetilde d_{t_-}(d)d^{-\sigma_{k,+}}
\inf_{a\in I_a}\left(1-|\cos(a\log d/2)|\right),
\quad
\sigma_{k,+}=(1+h)/2+t_+L_{k,+}/4.
\]
The higher upper abscissa supplies a coefficient lower bound valid at
both points. Therefore
\begin{equation}\label{eq:logpair}
\frac{S(x-a,h,t)+S(x+a,h,t)}2
>\frac{F_j+F_k+G_{k,a}}2+G^{\rm Euler}_{jk,a}.
\end{equation}
The finite logarithmic estimate is certified on the first 320 time rows
and the five bounded frequency classes. The remaining rows and the unbounded
class use their complete cone/Euler estimates. The assembler enforces these
domains and does not extrapolate beyond them.

A separate grouped-head alternative retains pairs with both indices in
$[2,2048]$, disjoint from the constant/head pairs. Its signed coefficients
and each complete group phase interval are retained; the same identity
\eqref{eq:phase} supplies the deficit. It is an alternative to
\eqref{eq:logpair}. For each class pair, the assembler selects the largest valid lower bound.
It then takes the minimum over all eleven class pairs. This selection rule
avoids counting any shared contribution twice.

\section{Retuned physical source and exact selection}\label{sec:retuned}
The physical observable remains
\[
 S(x,h,t)=-\Ima\frac{H_t'}{H_t}(x+ih),\qquad
 \frac{S(x-rh,h,t)+S(x+rh,h,t)}2\ge\ell_{i,r}.
\]
Probe nonvanishing is established before division, as in the retained
source theorem \cite{source}. No old lower bound is interpolated in $h$.
The all-frequency source region, native-cutoff treatment, two reflected
branches and corrected Gaussian error argument are unchanged.

\subsection{Coverage and production scope}
The new frozen plan contains 608 closed time cells
$[i/4000,(i+1)/4000]$, $0\le i<608$, through time $0.152$.
There are six frequency ranges per cell, including an unbounded final
range. Every same-class pair and every adjacent-class pair is included:
six plus five possibilities. The physical source starts at
$X_0-2$, where $X_0=6{,}000{,}000{,}185{,}827$ and $N_0=690{,}988$.
All future native cutoffs in the final range remain covered.

Relative to the old plan, exactly 274 heights change. The new heights
lie between $0.69$ and $1.88$; the old maximum was $2.18$.
For example, the first height becomes $47/25$ rather than $109/50$.
The source producer's enlarged analytic height domain already covers
this range. Its Cauchy-disk, activation, mass and denominator acceptance
conditions are checked at the actual time--height cells, not on an
invalid rectangle formed by pairing the smallest height with time zero.

All 608 cells are freshly produced in both constructions, including the
cells whose height did not change. Construction A uses 192-bit interval
arithmetic and integer-bin division 800; construction B uses 256 bits
and division 900, with a different derivative construction. Each is
produced in four adjoining blocks, with endpoints $0$--$151$,
$152$--$303$, $304$--$455$ and $456$--$607$. Complete arrays and receipts
are reconstructed after every block is present.
Both constructions use the same bundled Arb/FLINT library; they are not
independent transcendental implementations.

Exact source reconstruction checks 127,680 acceptance inequalities and
26,752 overlaps of the two constructions' temporal interval enclosures.
The corrected Gaussian and physical-error safety checks pass on all
1,216 error rows. The 256-bit and 384-bit finite-head wall files are
retained byte-for-byte. Those wall computations, and the published
finite-RH computation, are not represented as fresh production.

\subsection{Matched logarithmic coefficients at the new heights}
The joint translated-source producer is rerun in both constructions over
the complete schedule. Its exact reassembly checks 1,284,096 conditions.
The matched-logarithmic producer is then rerun at the same new heights,
with the original finite logarithmic head of 8,192 on the first 320
cells, common native head 690,988, and head-pair groups through 2,048 on
all 608 cells \cite{matchedsource}.

The improvement retains the original order of operations: common finite
convolution terms are cancelled before absolute values are taken, and a
complete bound is formed only after all remaining errors and the
logarithmic-derivative denominator are paid. This yields 9,120 positive
head-pair credits and 6,384 improvements from the matched-log alternative
in this reassembly. Cross-construction checks agree on 608 heat-coefficient
grids, 315 available logarithmic-phase grids, and 560 head-group phase
grids. A cache entry without a logarithmic grid is not treated as though
it had one.

The focused 16-band source producer is \emph{not} rerun at new heights.
Instead, the new valid broad continuum bound is used in every narrow
band. At cells with exactly the same height and time as the old source,
the maximum of the two complete valid bounds is allowed. This is a
maximum of whole certified estimates, never an unverified mixture of a
numerator from one construction and a denominator from another.

\subsection{A frozen whole-cell choice}
A rational discovery pass selects either the old enhanced source or the
new complete source in each cell. The resulting list of 608 provider
indices is frozen before the final histories are regenerated. It selects
466 new-source cells and 142 old-source cells; 274 selected heights
actually differ from the original heights. A new provider may improve
the floor even where its height is unchanged.

For every selected cell, the height and entire row of lower bounds are
copied from the chosen provider. The independent reconstruction checks
this operation verbatim and binds the resulting source to all four
comparison configurations by SHA-256. Neither a sampled optimum nor
greedy global optimality is a premise. Any complete, valid frozen choice
would give a legitimate comparison; this particular one gives
\eqref{eq:bound}.


\begin{proposition}[Selected physical source]\label{prop:source}
For every cell $I_i$ and either consumed shift $r=23/40$ or $r=3/5$,
the selected source supplies a rational lower bound $\ell_{i,r}$ such that
\[
 \frac{S(x-rh_i,h_i,t)+S(x+rh_i,h_i,t)}2>\ell_{i,r}
 \qquad(t\in I_i,\ x\ge X_0).
\]
The assertion includes all six frequency classes, all eleven translated
class pairs and every future native cutoff in the unbounded class.
\end{proposition}
\begin{proof}
The physical approximation and corrected remainder give
\eqref{eq:physical}. The common-phase estimate \eqref{eq:phase}, the
matched overlap \eqref{eq:resid} and the paired alternatives give valid
bounds for each class pair. The interval producers enclose their finite
masses, derivatives and phases on the complete time--height--frequency
cells. The exact reconstruction checks the activation and denominator
conditions, selects the strongest complete alternative for each pair,
and then takes the minimum over all eleven possibilities.
The selected source copies a complete valid row, including its height,
from one provider. The reconstruction checks this copy exactly.
Thus selection preserves each provider's all-frequency statement.
\end{proof}
\section{Full-domain kernel comparison}\label{sec:kernel}
At a simple highest upper zero $x+iy$, write $q=y^2$. The canonical
zero-motion identity, grouped into conjugate pairs, is
\begin{equation}\label{eq:motion}
 q'=-2-4y\sum_{\rm other\ pairs}K_y(d_j,b_j).
\end{equation}
The own conjugate contributes $-2$ exactly once. Real zeros have half-pair
weight, and the paired imaginary sums converge absolutely
\cite[Proposition 3.1]{polymath}. The task is to bound the remaining sum
from the source at the higher probes.

For $h>0$, define the conjugate-pair kernel
\begin{equation}\label{eq:kernel}
 K_h(d,b)=\frac{h-b}{d^2+(h-b)^2}
          +\frac{h+b}{d^2+(h+b)^2}.
\end{equation}
A translated observation at displacement $rh$ uses
\[
 J_{h,r}(d,b)=\frac{K_h(d+rh,b)+K_h(d-rh,b)}2.
\]
The useful inequality has the form
\begin{equation}\label{eq:comparison}
 2K_y(d,b)\ge\lambda\frac yh J_{h,r}(d,b),
 \qquad 0<y\le\rho h,\quad 0\le b\le y,\quad d\in\mathbb R.
\end{equation}
The point $(d,b)=(0,y)$ is excluded because $K_y$ is undefined there.
It is the distinguished contact pair, removed before this inequality is
summed. A multiple contact is treated by the separate contact argument;
no artificial value or path-independent limit is assigned at the singular
point.

\begin{proposition}[New certificate family]\label{prop:kernels}
All 44 triples $(r,\rho,\lambda)$ in Appendix~\ref{app:catalogue} satisfy
\eqref{eq:comparison}. The shifts are $r=23/40$ and $r=3/5$.
For each shift, the 22 values of $100\rho$ are
\[
 \begin{gathered}
 5,10,15,20,25,28,30,31,32,33,34,\\
 35,36,37,38,39,40,41,42,43,44,45.
 \end{gathered}
\]
Every $\lambda$ is an explicitly stored positive rational, at most two.
For example, the first two actively used domain refinements include
\[
 (r,\rho,\lambda)=\left(\frac{23}{40},\frac25,\frac{99641}{50000}\right),
 \qquad\left(\frac35,\frac25,\frac{198313}{100000}\right).
\]
\end{proposition}

\begin{proof}
Scale $h$ to one and put
\[
 D=d^2/h^2,\qquad v=y^2/h^2,\qquad B=b^2/h^2,\qquad R=r^2.
\]
Introduce the polynomials
\begin{align*}
 C&=D+R+1+B,& A&=C^2+4DR-4B,\\
 D_4&=A^2-16DRC^2,&G&=A(C-2B)-8DRC,\\
 Q&=(D+v-B)^2+4DB.
\end{align*}
Pairing the two translations first removes the square root of $D$.
After multiplication by positive denominators, \eqref{eq:comparison}
is equivalent to
\begin{equation}\label{eq:polynomial}
 P=(D+v-B)D_4-\frac{\lambda}{2}QG\ge0.
\end{equation}
The translated denominators are positive because $b/h\le\rho<1$.
The remaining denominator vanishes only at the excluded point.

Set $D=z/(1-z)$, $v=\rho^2p$, $B=\rho^2pq$. Then
\[
 \widetilde P(z,p,q)=(1-z)^5P\bigl(z/(1-z),\rho^2p,\rho^2pq\bigr)
\]
is a polynomial of degree at most five in each variable. The required
finite domain is the cube $0\le z,p,q\le1$, with $z=1$ included by
polynomial continuation.
Each certificate stores its rational coefficients, up to an independently
checked positive scalar, and a complete binary subdivision of this cube.
On every leaf, all tensor Bernstein coefficients are nonnegative. Since
the Bernstein basis is nonnegative on the cube, this proves the assertion.

The independent verifier does not import the searcher's polynomial algebra.
It evaluates the closed expression \eqref{eq:polynomial} and the stored
polynomial on the exact tensor grid
$\{0,1/6,\ldots,5/6\}^3$. The degree bound and six distinct nodes per
variable make this an unisolvent polynomial-identity check, not sampled
positivity. It also checks the positivity of the scalar normalisation.
For a degree-$n$ Bernstein coefficient vector $b$, the two half-interval
restrictions are independently computed by
\[
 b_k^L=2^{-k}\sum_{j=0}^{k}\binom kj b_j,\qquad
 b_k^R=2^{-(n-k)}\sum_{j=k}^{n}\binom{n-k}{j-k}b_j.
\]
These matrices are applied along the chosen coordinate at each split.
The partition verifier requires both children, forbids duplicate leaves
and forbids a leaf together with one of its descendants. Thus there is no
unverified gap in the cube cover.
The complete run checks 9,504 exact identities, 331 leaves and 71,496
Bernstein coefficients. Seven deliberately corrupted inputs are rejected.
\end{proof}

\subsection{The whole distinguished-pair subtraction}
At a zero of height $y$, write $w=y^2/h^2$. The distinguished pair's
normalised contribution at the translated probes is
\begin{equation}\label{eq:self}
 s_r(w)=\frac{2(1+r^2-w)}{(1+r^2-w)^2+4r^2w}.
\end{equation}
The constants used by the new certificates are
\[
 M_{23/40}=\frac{1503067131}{10^9},\qquad M_{3/5}=\frac{25}{17}.
\]
The bound $s_r(w)\le M$ for $0\le w\le\rho^2$ is verified by checking
nonnegativity of the quadratic
\begin{equation}\label{eq:selfpoly}
 Mw^2+\{M(2r^2-2)+2\}w+M(1+r^2)^2-2(1+r^2).
\end{equation}
Its endpoints and any interior minimum are checked exactly. The scalar
auditor separately checks the same condition on the full incoming
comparison interval at every step. An own-pair cap at just one sampled
height would not suffice.

If the physical translated source has lower bound $\ell$, summing
\eqref{eq:comparison} after removing the distinguished pair gives the rate
\begin{equation}\label{eq:rate}
 q'<-2-cq,\qquad
 c=\frac{2\lambda}{h}\left(\ell-\frac Mh\right)>0,
 \qquad q=y^2.
\end{equation}
The $-2$ is the own-conjugate contraction. The $M/h$ term subtracts the
entire contribution of that pair from the source; it is not optional.
Real zeros retain the half-weight convention of the zero-motion formula.


The history uses only fixed shifts. The retained code also supports a
continuum sharp mode, but no main-history step selects it. In that mode
one must evaluate the displacement at the actual contact height and
cover every swept displacement band; freezing the incoming displacement
with the sharp incoming self cap is invalid. The source supplement records
that correction explicitly. The present fixed-shift proof instead uses
\eqref{eq:selfpoly} on the entire incoming height interval.

\section{Global zero-height comparison}\label{sec:comparison}
Every time cell is divided into 32 equal pieces:
\[
\Delta t=\frac1{128000}.
\]
At each incoming certified upper value $q$, choose an admissible positive
rate \eqref{eq:rate}. Its comparison solution is
\begin{equation}\label{eq:flow}
Q(t+\tau)=\left(q+\frac2c\right)e^{-c\tau}-\frac2c.
\end{equation}
The program rounds each endpoint upwards and uses that upper bound to
start the next step. Another strip estimate or spatial guard never lowers
the restart value.

For a step starting at $q$, the selected certificate must satisfy
$q/h_i^2\le\rho^2$ and $|rh_i|<2$. The comparison decreases, so the same
kernel domain holds throughout the step. The own-pair polynomial is
checked on all $0\le w\le q/h_i^2$. These whole-interval tests are what
allow a fixed translation to be used safely; a self cap at just the
incoming height would not be enough.

\begin{lemma}[Global first-contact comparison]\label{lem:contact}
Under the hypotheses and certificates above, the piecewise flow
\eqref{eq:flow} and its upward restarts enclose the squared imaginary
height of every zero through the declared terminal step.
\end{lemma}
\begin{proof}
The first-rate check gives
\[
q_0-(2+c_0q_0)10^{-8}>9/16.
\]
The comparison is convex and above its tangent. Together with de Bruijn's
contraction from the initial strip, this establishes a positive interval
of strict separation.

On a closed time interval bounded away from zero, positive-time confinement
in \cite[Theorem 1.5(i)]{polymath} confines possible nonreal zeros to a
bounded real range. The wall localises a contact to $|x|>X_0$. Thus a
first-contact argument takes place on an actual compact zero set; no
possibly unattained global supremum is differentiated.

A multiple upper contact is impossible: the backward Hermite splitting of
\cite[Proposition 3.1(ii)]{polymath} produces an upper branch at an earlier
time with height excess of order $\sqrt\epsilon$, whereas the positive
barrier changes by $O(\epsilon)$. This contradicts previous separation.
At a simple contact all other zero heights are below the barrier, so the
full kernel inequality and source floor apply. The contact derivative is
strictly smaller than the comparison derivative by \eqref{eq:rate},
excluding a first crossing. Equivalently one can use the upper Dini
derivative of the finitely many local simple branches and integrating-factor
comparison. At a time knot the preceding left-hand inequality applies;
an upward restart cannot create a contact. These alternatives cover every
prospective first contact.
\end{proof}

Stop the final step at $q_e=10^{-8}$. The exact comparison travel time is
\begin{equation}\label{eq:lag}
\tau_e=\frac1c\log\left(1+\frac{c(q-q_e)}{2+cq_e}\right).
\end{equation}
The auditor verifies $0<\tau_e\le\Delta t$, ensuring that the last source
bound remains valid throughout the terminal interval. Lemma~\ref{lem:contact} and ordinary strip
contraction give
\begin{equation}\label{eq:terminal}
\Lambda\le t+\tau_e+\frac{q_e}{2}.
\end{equation}
The bound includes the full final strip-contraction time $q_e/2$.
If all zeros become real sooner, the same upper bound still holds.

\section{Exact endpoint and reproducibility}\label{sec:replay}
\subsection{Two rational constructions}
The primary producer bounds $e^{-x}$, $0\le x<1$, by its even degree-18
Taylor sum and rounds squared-height endpoints upward on the $2^{-96}$
grid. The alternate uses the reciprocal of the positive degree-12 Taylor
sum for $e^x$ and the $2^{-112}$ grid. For the final $\log(1+z)$,
$0<z<1$, the primary uses the odd degree-11 sum and the alternate uses
$z-z^2/(2(1+z))$. Both total terminal times are rounded upward to
multiples of $10^{-15}$.

The separate scalar auditor imports neither the producer nor its
selector. It uses an even degree-20 exponential sum and an odd degree-13
logarithmic sum with exact powers and factorials. It verifies the
hypothesis and seed, initial delay, source identity, every time join,
kernel eligibility, complete own-pair subtraction, every positive rate,
all outward endpoints and the terminal contraction. Agreement with the
producer is required at the exact inequality level, not just in displayed
decimal digits.

\begin{center}\small
\begin{tabular}{@{}lrr@{}}\toprule
Construction & Steps & Exact decimal ceiling\\\midrule
Primary & 14,899 & 0.116397726408291\\
Alternate & 14,899 & 0.116397726408312\\\bottomrule
\end{tabular}
\end{center}
Each construction uses $r=23/40$ on 1,024 steps and $r=3/5$ on 13,875
steps. The terminal source index is 465, and the final substep starts at
$7449/64000$. Its permitted source interval ends at $14899/128000$.
The exact final logarithmic travel time lies inside that interval.
Both histories have been regenerated byte-for-byte from the selected
source and checked by the separate auditor. Corrupted-history rejection
tests are retained with the verification receipts.

\begin{proof}[Proof of Theorem~\ref{thm:main}]
Lemma~\ref{lem:initial} gives the initial strip and strict separation.
The corrected wall localises every possible nonreal contact.
Proposition~\ref{prop:source} and the full-domain kernel certificates
supply the strict rate bound \eqref{eq:rate} at every simple contact;
Lemma~\ref{lem:contact} handles the global comparison, including multiple
contacts and time joins. The independently audited primary history
bounds the terminal travel time in \eqref{eq:lag}. Adding the complete
$q_e/2$ contraction allowance and rounding upward gives
\eqref{eq:bound}.
\end{proof}

\subsection{Separating the verification layers}
A hash identifies data; it does not establish a physical inequality.
The evidence therefore separates the analytic derivation, interval
production, exact reconstruction and scalar audit.

\begin{center}\small
\begin{tabularx}{\textwidth}{@{}p{31mm}XX@{}}\toprule
Obligation & Written argument & Executable check\\\midrule
Initial strip & Section~\ref{sec:norm}, conditional on $\HH_{7/8}$ & Seed and positive-delay tests\\[3pt]
Physical source & Sections~\ref{sec:physical}--\ref{sec:retuned}; source supplement; Foundations & Raw interval records, coverage and complete exact reconstruction\\[3pt]
Gaussian error and wall & Appendix~\ref{app:gaussian}; Foundations & Safety inequalities and retained wall certificates\\[3pt]
Kernel and self pair & Section~\ref{sec:kernel}; Appendix~\ref{app:catalogue} & Independent polynomial identity, full cube cover and self-cap tests\\[3pt]
Global comparison & Section~\ref{sec:comparison} & Complete joins and source binding\\[3pt]
Terminal bound & This section & Two rational histories and separate auditor\\\bottomrule
\end{tabularx}
\end{center}

\texttt{Proof\_and\_Replay.zip} contains the frozen scientific evidence,
raw source enclosures, original production receipts, analytic source
texts and all eight final-source histories. Only the two configurations
labelled \texttt{main} prove the headline under $\HH_{7/8}$; the other
six use separately stated initial strips. The \emph{Verification Guide}
gives the integrity-only, stored-source and native-production commands.
The stored-source audit reconstructs the full source and reruns the exact
comparisons. The optional native replay reruns the transcendental
producers at the frozen retuned plan. These are different operations.

The two native constructions use the same Arb/FLINT library
\cite{arb}, with different precision, integer partition and derivative
construction. The finite-head wall data are retained, not regenerated in
the retuning; the published finite-RH computation is not repeated.
The programs do not prove the external $7/8$ assertion or formalise the
complete analytic argument. No external referee certification,
unconditional record, global optimum or proof of RH is claimed.

\paragraph{Release and attribution.}
This is version 1.0 of the first deposit of the $0.116398$ conditional
paper. Earlier $0.118$ files remain in the immutable evidence as
\emph{undeposited development drafts}. The separate September record \cite{raymond149} retains its own
identity and review status.

\paragraph{AI assistance.}
The AI assisted with the derivations, code, numerical work and exposition.
The named author is responsible for the submitted claims. The supplied
checks and explicit assumptions, rather than this disclosure, define the
verification scope.

\appendix
\section{Exact kernel constants}\label{app:catalogue}
The table lists every new kernel certificate. Each displayed decimal is
exact. Each row has a complete rational polynomial and subdivision
certificate in \texttt{evidence/GEOMETRY\_CERTIFICATES.json}. The admissible
domain and its independent verification are proved in
Section~\ref{sec:kernel}. The whole-interval self caps are
$1503067131/10^9$ for $r=23/40$ and $25/17$ for $r=3/5$.

\begin{center}\small
\begin{tabular}{@{}rrr@{}}\toprule
$\rho$ & $\lambda$ for $r=23/40$ & $\lambda$ for $r=3/5$\\\midrule
0.05 & 1.99998 & 1.99819\\
0.1 & 1.99997 & 1.99784\\
0.15 & 1.99992 & 1.99719\\
0.2 & 1.99968 & 1.99612\\
0.25 & 1.99911 & 1.99448\\
0.28 & 1.99853 & 1.99312\\
0.3 & 1.998 & 1.99202\\
0.31 & 1.99768 & 1.9914\\
0.32 & 1.99733 & 1.99074\\
0.33 & 1.99694 & 1.99002\\
0.34 & 1.99651 & 1.98925\\
0.35 & 1.99603 & 1.98841\\
0.36 & 1.99551 & 1.98751\\
0.37 & 1.99493 & 1.98654\\
0.38 & 1.99429 & 1.98549\\
0.39 & 1.99359 & 1.98436\\
0.4 & 1.99282 & 1.98313\\
0.41 & 1.99197 & 1.98179\\
0.42 & 1.99105 & 1.98033\\
0.43 & 1.99003 & 1.97036\\
0.44 & 1.96906 & 1.93243\\
0.45 & 1.92508 & 1.89014\\
\bottomrule\end{tabular}
\end{center}
\section{Gaussian remainder correction and a safe wall}\label{app:gaussian}
This appendix replaces the $5/6$ Gaussian majorant in the retained
Foundations, Part A, Section 4.1, and the narrower historical
\nolinkurl{physical_eC0} wall comparison. The original archival source is
preserved. The broad physical error allowances used by both the original and retuned
source constructions remain valid.

Let $T'=T+\pi t/8\ge T$ and $D=T-3$. In the all-real-abscissa
Stirling/Taylor comparison accompanying \cite[(62)]{polymath}, the valid
coarse inequality is
\[
\frac1{12(T'/2-0.33)}\le\frac1{6(T-3)},
\]
not the expression with numerator $1/3$ over $4(T-3)$. With the corrected
numerator $2/3$,
\[
\frac{u^2+6|u|+2/3}{4D}\le\frac{u^2+11/12}{D},
\]
because multiplication of the difference by $4D$ gives
$3(|u|-1)^2$. The Gaussian integral is therefore
\begin{align}\label{eq:Gaussian}
\widetilde G_\sigma(c)
&=\int_{\mathbb R}e^{(u^2+11/12)/D+cu}
\frac{e^{-(u-\sigma)^2/t}}{\sqrt{\pi t}}\,\dd u\nonumber\\
&=\frac{\exp\left\{\dfrac{11}{12D}+
\dfrac{\sigma^2/D+\sigma c+tc^2/4}{1-t/D}\right\}}
{\sqrt{1-t/D}}.
\end{align}
At $t=0$ use its limiting value. Every actual disk satisfies
$-1\le\sigma\le2$, $0\le t\le1/5$ and
\[
D\ge(X_0-4)/2-3>10^{12}.
\]
The elementary inequality $-\log(1-z)\le z/(1-z)$ gives
\[
\log\widetilde G_\sigma(0)<6/D,
\quad\widetilde G_\sigma(0)-1<10^{-10}.
\]
For $|c|\le11/5$ the same explicit formula gives
$\log\widetilde G_\sigma(c)<93/20$ and
$\widetilde G_\sigma(c)<110$. This covers the required slopes
$\log9$ and $-\log2$. The numerical constant inequalities are enclosed
by rational positive Taylor sums and geometric tails.

The two all-real-$\sigma$ coefficient errors in
\cite[(62)--(64)]{polymath} therefore satisfy
\[
\delta_1+\delta_2
\le\frac{(0.2+0.029)110}{N_0-1}<10^{-3}.
\]
For the far negative Gaussian tail $u=-r$, $r\ge100$, the complete
coefficient majorant in the Foundations bounds the integrand by
\[
\exp\{-4r^2+20r+13+(r^2+1)/D\}\le e^{-3r^2}.
\]
The quadratic difference is positive and increasing from $r=100$ for
$D\ge100$, and
$\int_{100}^\infty e^{-3r^2}\dd r\le e^{-30000}/600<10^{-100}$.
Thus the full remainder bracket is bounded by
\[
\tfrac12(1+10^{-10})+
\frac{(0.2+0.029)110}{N_0-1}+10^{-100}<1.
\]
This establishes the retained coarse remainder bound:
\[
|R_{t,N}(\sigma+iT)|<e^{t\pi^2/64}|M_0(iT')|.
\]
Comparison with the physical normaliser then gives the same broad
physical error bound
\begin{equation}\label{eq:broad}
e_C^{\rm broad}
<4\exp\left\{-\frac{(1+y)L}{4}-\frac{tL^2}{16}+\frac1{10}\right\}.
\end{equation}
The existing producers subtract the allowance \eqref{eq:broad} on their
actual Cauchy disks, together with the complete edge pair and the
leading-model value and derivative error bounds. The separate safety checker verifies the 29 rational conditions and
1,216 physical-error rows for each complete two-construction source.
It is applied separately to the original and the retuned data.

At the wall \eqref{eq:broad} is largest at $t=y=0$. Using
$\pi<22/7$, $e^{1/10}<10/9$ and $7X_0/88>820^4$ gives
\[
e_C^{\rm broad}<\frac{40}{9\cdot820}=\frac2{369},
\qquad\frac2{369}+10^{-6}<\frac1{180}.
\]
The two wall certificates give ideal-model floor $79/100$, full ideal
displacement below $1/100$, and $e_A+e_B<10^{-6}$. Hence
\[
|H_t/B_t|>79/100-1/100-1/180=697/900>3/4.
\]
The narrower historical wall-error field is not used as a premise. The
archived documents and hashes remain unchanged, and the new receipt records
the replacement inequalities explicitly.

\begin{thebibliography}{12}
\setlength{\itemsep}{4pt}
\bibitem{raymond149} D. J. Raymond,
\emph{Correlated translated probes and the bound $\Lambda<0.149396$ for
the de Bruijn--Newman constant}, Zenodo preprint, 21 September 2026,
\href{https://doi.org/10.5281/zenodo.22867547}{doi:10.5281/zenodo.22867547}.
Includes \emph{Analytic Foundations} and \nolinkurl{Proof_and_Replay.zip}.
Public companion revision
\nolinkurl{3e519044ed8f2f969d76dc0c34008727302e67ec}.

\bibitem{openai78} OpenAI,
\emph{The Quasi-Riemann Hypothesis: A Zero-Free Half-Plane
$\Rea(s)>7/8$}, OpenAI Math Release preprint
\texttt{OAI:The-Quasi-Riemann-Hypothesis-September-30-2026}, 2026.
\href{https://github.com/openai/math/tree/fd4aeeb2ee4fc729c18d98444fed42fd0529eeeb/preprints/The-Quasi-Riemann-Hypothesis-September-30-2026}{Pinned manuscript directory}.

\bibitem{openai1112} OpenAI,
\emph{The Quasi-Riemann Hypothesis}, 5 October 2026.
\href{https://github.com/openai/math/tree/fd4aeeb2ee4fc729c18d98444fed42fd0529eeeb/preprints/The-Quasi-Riemann-Hypothesis-October-5-2026}{Pinned $11/12$ companion}.

\bibitem{openaiLean} OpenAI,
\nolinkurl{OAI.riemannZeta_ne_zero_of_seven_eighths_lt_re},
\href{https://github.com/openai/math/blob/fd4aeeb2ee4fc729c18d98444fed42fd0529eeeb/lean/OAI/NumberTheory/DirichletL/Nonvanishing.lean}{\texttt{Nonvanishing.lean}},
commit \nolinkurl{fd4aeeb2ee4fc729c18d98444fed42fd0529eeeb},
Pinned supplied source identity. No Lean execution is reported here.

\bibitem{polymath} D. H. J. Polymath,
\emph{Effective approximation of heat flow evolution of the Riemann
$\xi$ function, and a new upper bound for the de Bruijn--Newman constant},
Research in the Mathematical Sciences \textbf{6} (2019), article 31.
\href{https://doi.org/10.1007/s40687-019-0193-1}{doi:10.1007/s40687-019-0193-1}.
Numbered references follow \href{https://arxiv.org/abs/1904.12438v2}{arXiv:1904.12438v2}.

\bibitem{pt} D. J. Platt and T. S. Trudgian,
\emph{The Riemann hypothesis is true up to $3\cdot10^{12}$},
Bulletin of the London Mathematical Society \textbf{53} (2021), 792--797.
\href{https://doi.org/10.1112/blms.12460}{doi:10.1112/blms.12460};
\href{https://arxiv.org/abs/2004.09765}{arXiv:2004.09765}.

\bibitem{gomila} J. Gomila,
\emph{A computer-assisted proof of the bound
$\Lambda\le893927/5000000=0.1787854$ for the de Bruijn--Newman constant},
2026 manuscript cited in the retained earlier source.
\href{https://github.com/judegomila/dbn-lambda-01787854-candidate-audit/tree/a74738deb6d5e0f76887cb36901da08b68dca705}{Pinned source cited by the prior publication}; context only.

\bibitem{gordon} S. Gordon,
\emph{A computer-assisted upper bound for the de Bruijn--Newman constant},
2026 manuscript cited in the retained earlier source.
\href{https://github.com/stefangordon/dbn-upper-bound/tree/64c47f2db67288beb62c24d7a0f7aedd645dc4b3}{Pinned source cited by the prior publication}; context only.

\bibitem{planat} M. Planat,
\emph{Collective Contraction in the de Bruijn--Newman Heat Flow:
Off-Zero Logarithmic Derivatives and Certified Barrier Refinements},
2026, \href{https://arxiv.org/abs/2609.37164}{arXiv:2609.37164}.

\bibitem{arb} F. Johansson,
\emph{Arb: efficient arbitrary-precision midpoint-radius interval arithmetic},
IEEE Transactions on Computers \textbf{66} (2017), 1281--1292.
\href{https://doi.org/10.1109/TC.2017.2690633}{doi:10.1109/TC.2017.2690633}.
\bibitem{matchedsource} D. J. Raymond,
\emph{Matched logarithmic source: analytic supplement}, version 1.0,
10 October 2026. Included as \texttt{Matched\_Source\_Supplement.pdf};
original analytic source and producer code retained in the accompanying
\texttt{Proof\_and\_Replay.zip}.
\bibitem{source} D. J. Raymond, \emph{Coupled probe motion source},
\nolinkurl{BCLD_COUPLED_PROBE_MOTION_SOURCE_20260921},
\texttt{PROOF.md}, Sections 4--12. Original analytic text and producers
retained in the evidence dependency archive; typeset in
\emph{Prior Analytic Foundations}. The Gaussian replacement is
Appendix~\ref{app:gaussian} of the present paper.
\end{thebibliography}


\end{document}
